\subsection{Fisher-equivalent bandwidth and serial timing stages}
\label{sec:bfi}

For square-integrable conditional delay densities, Eq.~\eqref{eq:parseval} permits an exact equivalent-rectangular interpretation of the collision resource. Because $G$ is even and $G(0)=\eta$, define, for $\eta>0$, the ordinary-frequency DC-normalized Fisher-equivalent bandwidth
\begin{equation}
\boxed{
B_{\rm FI}
\equiv
\frac{1}{\eta}\int_0^\infty G(2\pi f)\,df.}
\label{eq:BFIdef}
\end{equation}
Changing variables from $f$ to $\omega$ and using Eq.~\eqref{eq:parseval} gives the exact identity
\begin{equation}
\boxed{
B_{\rm FI}=\frac{\Rtwo}{4\eta}.}
\label{eq:BFI}
\end{equation}
This is an area bandwidth of the source-normalized Fisher transfer, not an electrical amplitude bandwidth. Combining Eqs.~\eqref{eq:BFI} and \eqref{eq:R2H} yields
\begin{equation}
\boxed{
B_{\rm FI}\le\frac{\Hcap}{4\eta}.}
\label{eq:BFIH}
\end{equation}
If every captured-event conditional hazard obeys the common ceiling $\Lambda(m)\le\Lambda$, then $\Hcap\le\eta\Lambda$ and
\begin{equation}
\boxed{B_{\rm FI}\le\frac{\Lambda}{4}.}
\label{eq:BFILambda}
\end{equation}
For the exponential delay $f(t)=\Lambda e^{-\Lambda t}$, equality holds in Eq.~\eqref{eq:BFILambda}. Its half-power frequency is instead $\Lambda/(2\pi)$, so $B_{\rm FI}$ is not a relabeling of the conventional $-3\,$dB bandwidth.

The band-subspace result Eq.~\eqref{eq:universalBandCondition} gives the inverse resource cost an additional operational meaning. If every admissible weak waveform supported in $|\omega|\le\Omega$ must retain at least an absolute Fisher fraction $q$, then $G(\omega)\ge q$ almost everywhere on that band. Hence Eq.~\eqref{eq:parseval} requires
\begin{equation}
\pi\Rtwo
\ge\int_{-\Omega}^{\Omega}G(\omega)d\omega
\ge2\Omega q.
\end{equation}
With $B=\Omega/(2\pi)$,
\begin{equation}
\boxed{
\Rtwo\ge4Bq,
\qquad
\Hcap\ge4Bq.}
\label{eq:universalBandCost}
\end{equation}
Thus the same coefficient previously obtained from a flat-average task is also the necessary timing-resource cost for guaranteeing a minimum Fisher retention over an entire band-limited waveform subspace.

A simple cascade law makes the operational meaning explicit. Consider independent unmarked autonomous delay-only stages. For stage $j$, let the capture probability be $\eta_j$ and the delay characteristic function be $H_j$, so
\begin{equation}
G_j(\omega)=\eta_j|H_j(\omega)|^2.
\end{equation}
For two serial independent stages, delays add and capture probabilities multiply; therefore
\begin{equation}
\boxed{
G_{12}(\omega)=G_1(\omega)G_2(\omega).}
\label{eq:cascade}
\end{equation}
Thus unresolved independent timing stages multiply source Fisher transfer. This scalar factorization is restricted to independent unmarked delay-only stages; general retained marks remain governed by Eq.~\eqref{eq:G}.

As an analytic architecture example, take $k$ serial independent exponential waiting stages of common rate $\lambda$ and total capture probability $\eta$. The total delay is Erlang,
\begin{equation}
f_k(t)=\frac{\lambda^k t^{k-1}e^{-\lambda t}}{(k-1)!},
\end{equation}
with
\begin{equation}
H_k(\omega)=\left(\frac{\lambda}{\lambda+i\omega}\right)^k
\end{equation}
and hence
\begin{equation}
\boxed{
G_k(\omega)
=\eta\left(\frac{\lambda^2}{\lambda^2+\omega^2}\right)^k.}
\label{eq:ErlangG}
\end{equation}
Direct integration of $f_k^2$ gives
\begin{equation}
\boxed{
\frac{\Rtwo}{\eta}
=\lambda\frac{(2k-2)!}{4^{k-1}[(k-1)!]^2},}
\label{eq:ErlangR2}
\end{equation}
and therefore
\begin{equation}
\boxed{
B_{\rm FI}
=\frac{\lambda}{4}
\frac{\binom{2k-2}{k-1}}{4^{k-1}}.}
\label{eq:ErlangBFI}
\end{equation}
Using the central-binomial asymptotic,
\begin{equation}
B_{\rm FI}
\sim\frac{\lambda}{4\sqrt{\pi(k-1)}}
\qquad (k\to\infty).
\label{eq:ErlangAsymptotic}
\end{equation}
Thus serial unresolved stochastic registration stages progressively consume equivalent temporal Fisher bandwidth even when every stage has the same bare microscopic rate.
